\documentclass[11pt,french]{article}

\usepackage{amsmath,amsfonts,graphicx}
\renewcommand{\baselinestretch}{1.1}

\usepackage{enumerate,multirow}
\usepackage{verbatim}
%\usepackage{eclbip}
%\usepackage{pst-node}
%\include{psfig.sty}
\usepackage{url}
\usepackage{fancyhdr}
%\usepackage{slashbox}
\usepackage[colorlinks=true]{hyperref}

%\usepackage{soul}

\usepackage{ulem}

\parindent=0cm
\textwidth 6.25in
\textheight 21.9 cm
\topmargin -1.4 cm
\oddsidemargin 0in
\evensidemargin 0in

%%%% debut macro anti-orphelins %%%%
%\widowpenalty=10000
%\clubpenalty=10000
%\raggedbottom
%%%% fin macro %%%%

\headsep 1.5cm
\footskip 1.5cm
\fancyhead{}
\fancyfoot{}
\renewcommand{\headrulewidth}{0.6pt}
%\renewcommand{\footrulewidth}{0.6pt}
\fancyhead[LO,LE]{{\footnotesize Université Laval \\
Faculté des sciences et de génie\\
Département de mathématiques et de statistique
}}
\fancyhead[RO,RE]{{\footnotesize Automne 2025\\STT-2902 Modélisation statistique\\Julien Miron}}
%\fancyfoot[LO, LE] {\includegraphics[width=1.5cm]{UL.jpg}}
%\fancyfoot[RO, RE] {\vspace{-0.5cm} \thepage}


\pagestyle{fancy}
\pagenumbering{gobble}


\begin{document}
\begin{center}
   \textbf{ Analyse de données -- Inférence}
\end{center}


\subsection*{Jeu de données}

Nous utiliserons le jeu de données \texttt{penguins} provenant du package \texttt{palmerpenguins} en R. Il contient des informations sur trois espèces de manchots vivant dans les îles Palmer, en Antarctique. Chaque ligne représente un manchot mesuré dans le cadre d'une étude écologique.

Variables utilisées :
\begin{itemize}
  \item \texttt{species} : Espèce du manchot (\textit{Adelie}, \textit{Chinstrap}, \textit{Gentoo})
  \item \texttt{flipper\_length\_mm} : Longueur des nageoires (en mm)
  \item \texttt{body\_mass\_g} : Masse corporelle (en grammes)
  \item \texttt{sex} : Sexe du manchot (\textit{male}, \textit{female})
  \item \texttt{island} : Île d’origine du manchot (\textit{Biscoe}, \textit{Dream}, \textit{Torgersen})
\end{itemize}

\subsection*{Questions de recherche}

\begin{enumerate}
  \item Est-ce que les manchots \textit{Gentoo} sont en moyenne plus lourds que les manchots \textit{Adelie} ?
  \item Est-ce que la proportion de manchots mâles est différente entre les îles \textit{Dream} et \textit{Biscoe}? Vous devez supposer que que les conditions sont bien remplies. Vous en reparlerez au point 3.
  \item Est-ce que la longueur des nageoires des manchots suit une distribution normale ?
\end{enumerate}

\subsection*{Votre travail}

Pour les questions 1 et 2 :

\begin{enumerate}
  \item \textbf{Statistiques descriptives} : Présentez des résumés numériques et graphiques. Essayez de faire une inférence informelle à partir du jeu de données.

  \item \textbf{Analyse inférentielle} : Utilisez un test d’hypothèse approprié (test t, test de proportion, test de normalité, etc.) avec un seuil de signification $\alpha = 0{,}01$. Fournissez aussi un intervalle de confiance à 90\% pour le paramètre étudié.

  \item \textbf{Postulats} : Vous avez pris certaines choses pour acquis au point 2. Quelles sont-elles? Si ces choses n'étaient pas vraies, discutez des limites de l’analyse.

  \item \textbf{Valeurs manquantes} : Si vous devez exclure certaines lignes en raison de valeurs manquantes, mentionnez-le clairement et discutez des impacts potentiels.
\end{enumerate}

Pour la question 3, utilisez des statistiques descriptives.

Vous pouvez effectuer les analyses dans R, Excel ou GeoGebra.
Incluez dans votre rapport des captures d’écran ou extraits de sorties de logiciel pour appuyer votre interprétation des résultats.

\end{document}